
\subsubsection{Analysis and the \lstinline[language=lean]|TopologicalSpace| Class}

We have briefly seen how one can build rational numbers from
natural numbers and integers using quotient types.
In our implementation of \lstinline[language=lean]|myRat|, we
invoked \lstinline[language=lean]|Classical| axioms merely for brevity,
to avoid writing an explicit decision procedure for the inequality.
However, when moving to real numbers, reliance on non-constructive tools
becomes much more unavoidable.
In Mathlib, real numbers are defined as equivalence classes of
Cauchy sequences of rational numbers \cite{mathlibdoc}.
Although experimental libraries for computable reals in Lean exist, such as the
exact real arithmetic implementation by Christiansen and Licata \cite{christiansen2022computing},
the standard \lstinline[language=lean]|Real| type used for topology relies
on classical axioms. Consequently, most operations on it are marked as
\lstinline[language=lean]|noncomputable|, indicating they cannot
be executed algorithmically.
For example, the sine function requires this marker:
\begin{lstlisting}[language=lean]
noncomputable def realSin (x : ℝ) : ℝ := Real.sin x
\end{lstlisting}
Real numbers are instances of
various algebraic and topological
structures like: linear ordered field, normed space, metric space, and
topological space (\cite{mathlib2020}).
In the next section, I will present an example of formalization that requires
working with real numbers their topological space.
The latter provides
foundational tools for analysis concepts like continuity and convergence.
Given that $\mathbb{R}$ is a metric space, these notions in can,
of course, be defined
using the familiar $\varepsilon$
formulation:
\begin{lstlisting}[language=lean]
def ConvergesTo (s : ℕ → ℝ) (a : ℝ) :=
  ∀ ε > 0, ∃ N, ∀ n ≥ N, |s n - a| < ε
\end{lstlisting}
where \lstinline[language=lean]|s| is a sequence
and \lstinline[language=lean]|a| is the limit point.
More generally, these concepts are formalized in topology.
In standard textbooks, a topological space
is defined as a set equipped with a collection of open sets
satisfying certain axioms. This is reflected in
Mathlib's \lstinline[language=lean]|TopologicalSpace| type class:
\begin{lstlisting}[language=lean]
class TopologicalSpace (α : Type*) where
  IsOpen : Set α → Prop
  isOpen_univ : IsOpen univ
  isOpen_inter : ∀ s t, IsOpen s → IsOpen t → IsOpen (s ∩ t)
  isOpen_sUnion : ∀ s, (∀ t ∈ s, IsOpen t) → IsOpen (sUnion s)
\end{lstlisting}
Recall that Lean treats sets as predicates
(\lstinline[language=lean]|Set α := α → Prop|),
where set membership is function application.
Here, \lstinline[language=lean]|IsOpen| is a predicate on sets indicating whether
they are open. The three axioms correspond to the standard topological axioms:
\begin{itemize}
    \item The whole space is open (\lstinline[language=lean]|isOpen_univ|)
    \item The intersection of two open sets is open (\lstinline[language=lean]|isOpen_inter|)
    \item The union of any collection of open sets is open (\lstinline[language=lean]|isOpen_sUnion|)
\end{itemize}
\begin{note}
    The empty set is open, as it is the union of an empty collection of open sets.
\end{note}
Using this structure, global continuity can be defined in terms
of open sets, and, indeed, this is how it is defined in Mathlib:
\begin{lstlisting}[language=lean]
structure Continuous (f : X → Y) : Prop where
  isOpen_preimage : ∀ s, IsOpen s → IsOpen (f ⁻¹' s)
\end{lstlisting}
However, when proving results about limits, convergence, and continuity,
Mathlib uses an alternative way based on \textbf{filters}.
This approach may seem more abstract initially, but it provides
a more general and powerful framework that works uniformly
across all topological spaces and, by extension, metric spaces.

\subsubsection{Limits and Convergence with Filters}

The concept of a limit is quiet extended, there are many
types of limits to consider.
For intance the limit of a sequence or a function at a point,
limits at infinity (from above or below), one-sided limits (from the left or right)
and so on.
Defining each of these separately would require a
huge amount of work to include in Mathlib,
with significant duplication of theorems and proofs.
Moreover, many fundamental theorems
(like the characterization of continuity via limits)
would need to be reproved for each type of limit.
Fortunately, Bourbaki solved this issue by introducing the notion
of \textbf{filters} to unify
all concepts of limits, convergence, neighborhoods and terms like eventually or
frequently often into a single framework.
Intuitively, a filter represents a notion of "sufficiently large" subsets.
More fomrally, a filter $F$ on a type $X$ is a collection of subsets
of $X$ satisfying three axioms:
\begin{lstlisting}[language=lean]
structure Filter (α : Type*) where
  sets : Set (Set α)
  univ_sets : Set.univ ∈ sets
  sets_of_superset {x y} : x ∈ sets → x ⊆ y → y ∈ sets
  inter_sets {x y} : x ∈ sets → y ∈ sets → x ∩ y ∈ sets
\end{lstlisting}
The field sets suggests to think of a filter as
a collection sets. The three axioms correspond to:
\begin{itemize}
    \item  $X \in F$ (the whole space is in the filter)
    \item  If $U \in F$ and $U \subseteq V$, then $V \in F$
          (supersets of "large" sets are "large")
    \item  If $U, V \in F$, then $U \cap V \in F$
          (finite intersections of "large" sets are "large")
\end{itemize}
Note that if $F$ is a filter that contains the empty set,
then it contains all subsets of $X$.
Filters are quiet categorical objects, and we are going
to intuivetly make use of them instead of descibing it fomrally.
In particular, we are going to use some of the following concetps:
\begin{itemize}
    \item \textbf{Neighborhood filter} \lstinline[language=lean]|𝓝 x|:
          In a topological space,
          this filter contains all neighborhoods of the point $x$.
          This recall the standard defintion of neighborhoods in topology.
          A set is in \lstinline[language=lean]|𝓝 x| if it contains an
          open set containing $x$. This captures the idea of "near $x$."
    \item \textbf{At top filter} \lstinline[language=lean]|atTop : Filter ℕ|: Contains sets that
          include all sufficiently large natural numbers.
          Formally, $U \in$ \lstinline[language=lean]|atTop| if and only if there
          exists $N$ such that $\{n \mid n \geq N\} \subseteq U$.
          This captures the idea of "$n \to \infty$."
    \item \lstinline[language=lean]|∀ᶠ x in f, p x| (\lstinline[language=lean]|f.Eventually p|):
          "Eventually in filter $f$, property $p$ holds."
          This means there exists some set $U \in f$ such that $p$ holds for all $x \in U$.
          % For example, \lstinline[language=lean]|∀ᶠ n in atTop, n > 100| means
          % "for all sufficiently large $n$, we have $n > 100$."
    \item \lstinline[language=lean]|∃ᶠ x in f, p x| (\lstinline[language=lean]|f.Frequently p|):
          "Frequently in filter $f$, property $p$ holds." This means for every set $U \in f$, there exists
          some $x \in U$ where $p$ holds. This captures the idea that $p$ holds "infinitely often" or
          "arbitrarily close."
          % For example, \lstinline[language=lean]|∃ᶠ n in atTop, Even n| means
          % "there are arbitrarily large even numbers."
    \item \lstinline[language=lean]|Tendsto f l₁ l₂|: "Function $f$ tends from filter
          $l_1$ to filter $l_2$." This is used for convergence.
\end{itemize}
\begin{example}
    We can express convergence of a sequence $s_n$ to its limit $a$ using filters.
    \lstinline[language=lean]|Tendsto s atTop (𝓝 a)|; meaning $s_n \to a$ as $n \to \infty$
\end{example}
\begin{example}
    Another more insigthfull example is the definition of local continuity;
    at a point $x$ or restricted to a subset.
    As seen before, the structure \lstinline[language=lean]|Continuous|,
    for global continuity, is defined
    in terms of open sets. However, Mathlib
    also provides alternative definitions
    for local continuity: \lstinline[language=lean]|ContinuousAt|
    defines continuity at a single point,
    \lstinline[language=lean]|ContinuousWithinAt| defines continuity within a
    set at a point, and
    \lstinline[language=lean]|ContinuousOn| defines continuity on an entire set.
    All these local characterizations are defined in terms of filters.
    The connection between the global continuity definition and the filter-based
    local definitions is established by the following fundamental theorem:
    \begin{lstlisting}[language=lean]
theorem Continuous.tendsto (hf : Continuous f) (x) : 
    Tendsto f (𝓝 x) (𝓝 (f x)) :=
  ((nhds_basis_opens x).tendsto_iff <| nhds_basis_opens <| f x).2 fun t ⟨hxt, ht⟩ =>
    ⟨f ⁻¹' t, ⟨hxt, ht.preimage hf⟩, Subset.rfl⟩
\end{lstlisting}
    The key concept here is \lstinline[language=lean]|FilterBasis|.
    Similar to how a topological space can be defined via a basis of open sets,
    a filter can be defined via a basis of sets.
    A basis  $B$ for a filter $F$ is a nonemprty collection of sets wich
    preserve intersections;
    for any two sets $U, V \in B$, there exists a set $W \in B$ such
    that $W \subseteq U \cap V$.
    \begin{lstlisting}[language=lean]
  structure FilterBasis (α : Type*) where
    sets : Set (Set α)
    nonempty : sets.Nonempty
    inter_sets {x y} : x ∈ sets → y ∈ sets → ∃ z ∈ sets, z ⊆ x ∩ y
  \end{lstlisting}
    Given a basis, we can not only generate a filter,
    but also help proving properties about it, restricitng to one basis only.
    \newpage
    To better understand the proof, we can expand it slightly:
    \begin{lstlisting}[language=lean]
theorem Continuous'.tendsto (hf : Continuous f) (x) :
    Tendsto f (𝓝 x) (𝓝 (f x)) := by
  rw [(nhds_basis_opens x).tendsto_iff (nhds_basis_opens (f x))]
  intro t ⟨hft_in, ht_open⟩
  use f ⁻¹' t
  constructor
  · exact ⟨hft_in, ht_open.preimage hf⟩
  · exact Subset.rfl
\end{lstlisting}
    The statemnt \lstinline[language=lean]|Tendsto f (𝓝 x) (𝓝 (f x))| is reformulation of
    continuity in terms of neighborhood filters \lstinline[language=lean]|𝓝 x|.
    Meaning that for every neighborhood of $f(x)$, there exists a neighborhood of $x$
    that maps into it under $f$.
    In particular the neighborhood filter \lstinline[language=lean]|𝓝 x| has a basis
    consisting of all open sets containing the point $x$,
    wich is stated by \lstinline[language=lean]|nhds_basis_opens x|.
    Thus the line
    \lstinline[language=lean]|rw [(nhds_basis_opens x).tendsto_iff (nhds_basis_opens (f x))]|
    restricts the goal in terms of these basis open sets.
    This converts our goal to: for every open neighborhood (set) $t$ of $f(x)$,
    there exists an open neighborhood $s$ of $x$ with $f^{-1} s \subseteq t$.
    Next, we introduce an arbitrary open neighborhood $t$ of $f(x)$
    using \lstinline[language=lean]|intro t ⟨hft_in, ht_open⟩|.
    The natural choice for the witness neighborhood $s$ is
    given by \lstinline[language=lean]|use f ⁻¹' t|.
    Now, the first part of the goal is to show that $x \in f^{-1}(t)$
    and that $f^{-1}(t)$ is open.
    This is done using
    \lstinline[language=lean]|constructor; exact ⟨hft_in, ht_open.preimage hf⟩|.
    With \lstinline[language=lean]|ht_open.preimage hf| we are
    leveraging the global continuity of $f$. Since $t$ is open and $f$ is continuous,
    the preimage $f^{-1}(t)$ is also open.
    The second part of the goal is to show that $\forall x \in f^{-1}(t), f(x) \in t$,
    which always holds by the definition of preimage; \lstinline[language=lean]|exact Subset.rfl|.
    Here the full example: [\href{https://live.lean-lang.org/#codez=JYWwDg9gTgLgBAWQIYwBYBtgCMBQOJgCmAdnACoEToQDmAnnAGLDoyFRwDKhMOAbkijAkWdITgBvABpwAmnABacAFzk6RAFQBffoOGjxAbQqRqNYAGMk6TmCQXxUgLpxjlM5eu3742U91CImJwABQAZipwMoBJhHIAlKEAHpFScXhohNCEIHAAwhDEMMDEAK4QJQDOAOQ1AHRsxAAmFTAQoagRqvmFxWWVcGEJIYkJyjhw5CTNrQOhgKwbgLs7cCPzS+HLcaMAvHBYdONwUADuriHEqM0A+lhIFcAVlwQkFRv1Uy0Ql8BhEWcXDzc7g8nsQXusRnF/BNijAoG14IAL8g6MC+xAANHBUCiQYBL8gOlXEEUA3gQATqqcF4EwsBRaUBKFlaUAOAHa4IREvZEcjURisY8iMRamAoIRQEgaOIOniJqz2ZyuCUsBUeLUoGF0EA}{link to Lean live}]

    Using this bridge from global to local continuity,
    we can understand the following connection:
    \begin{lstlisting}[language=lean]
def ContinuousAt (f : X → Y) (x : X) :=
  Tendsto f (𝓝 x) (𝓝 (f x))
theorem Continuous.continuousAt (h : Continuous f) : ContinuousAt f x :=
  h.tendsto x
\end{lstlisting}
    The remaining local continuity concepts are defined similarly:
    \begin{lstlisting}[language=lean]
def ContinuousWithinAt (f : X → Y) (s : Set X) (x : X) : Prop :=
  Tendsto f (𝓝[s] x) (𝓝 (f x))
def ContinuousOn (f : X → Y) (s : Set X) : Prop :=
  ∀ x ∈ s, ContinuousWithinAt f s x
\end{lstlisting}
    Note that \lstinline[language=lean]|𝓝[s] x| denotes the neighborhood filter
    of $x$ restricted
    to the set $s$, allowing us to study the behavior of functions on arbitrary subsets.
    Next we will use \lstinline[language=lean]|ContinuousOn| in our formalization.
\end{example}
